%% kapa: exam=1 round=34 year=2023 period=2 form=A subject=appraisal_law subject_ko=감정평가관계법규 seq=22 number=22
%% source: Q-Net 감정평가사 기출문제 (cst00309) artlSeq=5211789 file=raw/1008/5211789/2023년도 제34회 감정평가사 1차시험 2교시.pdf
%% conversion: kordoc 4.24.51 -> tools/gen1.py (auto); review=proofread-sonnet; flags=none
\begin{question}{22}
국유재산법령상 행정재산의 사용허가와 일반재산의 처분에 있어 두 번의 일반경쟁입찰에도 낙찰자가 없는 경우 세 번째 입찰부터의 예정가격 조정·결정의 방법에 관한 설명이다. (　　)에 들어갈 숫자로 옳은 것은?
\begin{framedtext}
○ 행정재산 사용허가의 경우: 최초 사용료 예정가격의 100분의 (ㄱ)을 최저한도로 하여 매회 100분의 10의 금액만큼 그 예정가격을 낮추는 방법으로 조정할 수 있다. \par
○ 일반재산 처분의 경우: 최초 매각 예정가격의 100분의 (ㄴ)을 최저한도로 하여 매회 100분의 10의 금액만큼 그 예정가격을 낮출 수 있다. \par
\end{framedtext}
\begin{choices}
  \item ㄱ: 10, ㄴ: 30
  \item ㄱ: 10, ㄴ: 50
  \item ㄱ: 20, ㄴ: 30
  \item ㄱ: 20, ㄴ: 50
  \item ㄱ: 30, ㄴ: 50
\end{choices}
\end{question}
\answer{4}
